\documentclass[10pt,a4paper]{article}
\usepackage[utf8]{inputenc}
\usepackage[vietnamese]{babel}
\usepackage{amsmath,amssymb,amsfonts}
\usepackage{geometry}
\geometry{a4paper, left=10mm, right=10mm, top=10mm, bottom=10mm}
\usepackage{tikz}
\usetikzlibrary{shapes,arrows.meta,positioning,calc}
\usepackage[most]{tcolorbox}
\usepackage{array,booktabs,multicol}
\usepackage{fontawesome}
\usepackage{helvet}
\renewcommand{\familydefault}{\sfdefault}
\usepackage{parskip}

\definecolor{primary}{RGB}{112, 48, 160}
\definecolor{teal}{RGB}{0, 128, 128}
\definecolor{orange}{RGB}{230, 108, 44}
\definecolor{green}{RGB}{46, 139, 87}
\definecolor{rose}{RGB}{192, 0, 0}

\pagestyle{empty}

\begin{document}

% --- HEADER BANNER ---
\begin{tcolorbox}[
  enhanced,
  colback=violet!8!white,
  colframe=violet!80!black,
  arc=2.5mm, boxrule=1pt,
  left=3mm, right=3mm, top=2mm, bottom=2mm
]
  \begin{tabular*}{\textwidth}{@{\extracolsep{\fill}} lr}
    {\small\bfseries\color{violet!85!black}\faBook\ SỔ TAY KIẾN THỨC TOÁN 2} & 
    {\small\bfseries\color{teal!85!black}\faGraduationCap\ CHƯƠNG TRÌNH GDPT 2018}
  \end{tabular*}
  \vspace{-1mm}
  \begin{center}
    \rule{0.98\textwidth}{0.6pt}
  \end{center}
  \vspace{-2mm}
  \centering
  {\LARGE\bfseries\color{violet!90!black}\faBookmark\ BÀI 51: SỐ CÓ BA CHỮ SỐ}\\[1mm]
  {\itshape\footnotesize\color{black!75} Cấu tạo số có 3 chữ số, cách đọc, viết theo hàng trăm, chục, đơn vị}
\end{tcolorbox}

\vspace{-2mm}

\noindent
\begin{minipage}[t]{0.49\textwidth}\vspace{0pt}
\begin{tcolorbox}[
  enhanced,
  colback=violet!4!white,
  colframe=violet!75!black,
  arc=2mm, boxrule=0.8pt,
  left=2mm, right=2mm, top=1.8mm, bottom=1.8mm,
  title={\footnotesize\bfseries\color{white}\faKey\ 1. Cấu tạo số có ba chữ số}
]
{\footnotesize
\begin{itemize}\setlength{\itemsep}{1pt}\setlength{\leftskip}{-2.5mm}
  \item Số có ba chữ số gồm 3 hàng (tính từ phải sang trái):
  \begin{itemize}\setlength{\itemsep}{0.5pt}\setlength{\leftskip}{-2mm}
    \item Hàng \textbf{đơn vị} (bên phải cùng).
    \item Hàng \textbf{chục} (ở giữa).
    \item Hàng \textbf{trăm} (bên trái cùng, chữ số khác $0$).
  \end{itemize}
  \item \textbf{Ví dụ:} Số $243$ gồm:
  \begin{itemize}\setlength{\itemsep}{0.5pt}\setlength{\leftskip}{-2mm}
    \item $2$ trăm, $4$ chục và $3$ đơn vị.
  \end{itemize}
  \item Số bé nhất có 3 chữ số là \textbf{$100$}.
  \item Số lớn nhất có 3 chữ số là \textbf{$999$}.
\end{itemize}
}
\end{tcolorbox}
\end{minipage}
\hfill
\begin{minipage}[t]{0.49\textwidth}\vspace{0pt}
\begin{tcolorbox}[
  enhanced,
  colback=teal!4!white,
  colframe=teal!75!black,
  arc=2mm, boxrule=0.8pt,
  left=2mm, right=2mm, top=1.8mm, bottom=1.8mm,
  title={\footnotesize\bfseries\color{white}\faSliders\ 2. Cách đọc và viết số có ba chữ số}
]
{\footnotesize
\begin{itemize}\setlength{\itemsep}{1pt}\setlength{\leftskip}{-2.5mm}
  \item \textbf{Thứ tự:} Đọc và viết từ trái sang phải, từ hàng trăm đến hàng chục rồi đến hàng đơn vị.
  \item \textbf{Trường hợp đặc biệt ở hàng chục và đơn vị:}
  \begin{itemize}\setlength{\itemsep}{0.5pt}\setlength{\leftskip}{-2mm}
    \item Chữ số hàng chục là $0$: đọc là \textbf{"linh"} (ví dụ: $305$ đọc là \textit{Ba trăm linh năm}).
    \item Chữ số hàng chục $\ge 2$ và hàng đơn vị là $1$: đọc là \textbf{"mốt"} (ví dụ: $231$ đọc là \textit{Hai trăm ba mươi mốt}).
    \item Hàng đơn vị là $5$ và hàng chục $\ge 1$: đọc là \textbf{"lăm"} (ví dụ: $415$ đọc là \textit{Bốn trăm mười lăm}).
  \end{itemize}
\end{itemize}
}
\end{tcolorbox}
\end{minipage}

\vspace{-1.5mm}
\begin{tcolorbox}[
  enhanced,
  colback=orange!4!white,
  colframe=orange!80!black,
  arc=2mm, boxrule=0.8pt,
  left=2mm, right=2mm, top=1.5mm, bottom=1.5mm,
  title={\footnotesize\bfseries\color{white}\faSitemap\ 3. Sơ đồ cấu tạo bảng hàng số có 3 chữ số}
]
\centering
\begin{tikzpicture}[x=1.6cm, y=0.75cm]
  \tikzset{arr/.style={-{Stealth[scale=0.8]}}, every node/.append style={align=center}}
  \node[rectangle, rounded corners=2pt, draw=violet!80!black, fill=violet!15, font=\scriptsize\bfseries] at (0,0) {HÀNG TRĂM\\(2 tấm thẻ 100)};
  \node[rectangle, rounded corners=2pt, draw=teal!80!black, fill=teal!15, font=\scriptsize\bfseries] at (2.2,0) {HÀNG CHỤC\\(4 thanh 1 chục)};
  \node[rectangle, rounded corners=2pt, draw=orange!80!black, fill=orange!15, font=\scriptsize\bfseries] at (4.4,0) {HÀNG ĐƠN VỊ\\(3 khối lập phương)};

  \node[rectangle, rounded corners=2pt, draw=green!70!black, fill=green!15, font=\footnotesize\bfseries] at (2.2,-1.2) {
    Viết số: $\mathbf{243}$ \quad --- \quad Đọc số: \textit{Hai trăm bốn mươi ba}
  };
  \draw[arr, thick] (0,-0.4) |- (2.2,-0.9);
  \draw[arr, thick] (2.2,-0.4) -- (2.2,-0.9);
  \draw[arr, thick] (4.4,-0.4) |- (2.2,-0.9);
\end{tikzpicture}
\end{tcolorbox}

\vspace{-1.5mm}
\noindent
\begin{minipage}[t]{0.49\textwidth}\vspace{0pt}
\begin{tcolorbox}[
  enhanced,
  colback=green!4!white,
  colframe=green!60!black,
  arc=2mm, boxrule=0.8pt,
  left=2mm, right=2mm, top=1.8mm, bottom=1.8mm,
  title={\footnotesize\bfseries\color{white}\faCheckCircle\ 4. Ví dụ mẫu chuẩn mực}
]
{\footnotesize
\textbf{1. Đọc các số sau:}\\
- $356$: Ba trăm năm mươi sáu.\\
- $508$: Năm trăm linh tám.\\
- $615$: Sáu trăm mười lăm.\\
- $721$: Bảy trăm hai mươi mốt.\\
\textbf{2. Viết số:}\\
- Bốn trăm linh hai: $\mathbf{402}$.\\
- Chín trăm chín mươi: $\mathbf{990}$.\\
- Tám trăm bốn mươi lăm: $\mathbf{845}$.
}
\end{tcolorbox}
\end{minipage}
\hfill
\begin{minipage}[t]{0.49\textwidth}\vspace{0pt}
\begin{tcolorbox}[
  enhanced,
  colback=rose!4!white,
  colframe=rose!70!black,
  arc=2mm, boxrule=0.8pt,
  left=2mm, right=2mm, top=1.8mm, bottom=1.8mm,
  title={\footnotesize\bfseries\color{white}\faExclamationTriangle\ 5. Cạm bẫy \& Lưu ý chính tả}
]
{\footnotesize
\begin{itemize}\setlength{\itemsep}{1.5pt}\setlength{\leftskip}{-2.5mm}
  \item \textbf{Phân biệt "mười" và "mươi":}\\
  $15$ là \textit{mười lăm}, nhưng $50$ là \textit{năm mươi}.
  \item \textbf{Phân biệt "một" và "mốt":}\\
  $101$: Một trăm linh \textbf{một}.\\
  $121$: Một trăm hai mươi \textbf{mốt}.
  \item \textbf{Phân biệt "năm" và "lăm":}\\
  $305$: Ba trăm linh \textbf{năm}.\\
  $325$: Ba trăm hai mươi \textbf{lăm}.
\end{itemize}
}
\end{tcolorbox}
\end{minipage}


\end{document}
